\documentclass[a4paper, 12pt]{article}

% Packages
\usepackage[utf8]{inputenc}
\usepackage{amsmath, amssymb, amsthm}
\usepackage{graphicx}
\usepackage{hyperref}
\usepackage[margin=2.5cm]{geometry}

% Document information
\title{Logica, verzamelingen en relaties \textemdash{} Huiswerk 1}
\author{Harman Preet Singh}

% Custom theorem environments
\newtheorem{theorem}{Theorem}
\newtheorem{lemma}[theorem]{Lemma}
\newtheorem{proposition}[theorem]{Proposition}
\newtheorem{corollary}[theorem]{Corollary}
\newtheorem{definition}{Definition}
\newtheorem{example}{Example}

\begin{document}
% chktex-file 44

\maketitle
% \tableofcontents

\section{Opgave 1}
\textbf{A}

\[
\begin{tabular}{|c|c|c|c|c|}
	\hline
	${p}$ & ${q}$ & ${q}\rightarrow{p}$ & $\neg{p}$ & $\neg{p}\lor({q}\rightarrow{q})$ \\
	\hline
	1 & 0 & 1 & 0 & 1 \\
	1 & 1 & 1 & 0 & 1 \\
	0 & 1 & 0 & 1 & 1 \\
	0 & 0 & 1 & 1 & 1 \\
	\hline
\end{tabular}
\]

In alle vier gevallen is $\phi$ waar, dus het is een tautologie.

\vspace{2cm}

\noindent\textbf{B}

Ik gebruik eerst de standaardequivalentie $q \rightarrow p$ = $\neg q \lor p$. Dus:

\begin{center}
	$\phi$ = $\neg p \lor (q \rightarrow p)$ = $\neg p \lor (\neg q \lor p)$ \\
\end{center}

Dan verander ik de haakjes en de volgorde van de termen (mag door de commutatieve en associatieve wetten voor $\lor$):
\begin{center}
	$\phi$ = $(\neg p \lor p) \lor \neg q$ \\
\end{center}

En dan pas ik de wet van uitgesloten derde toe op ($\neg p \lor p$). Dus ($\neg p \lor p$) schrijf ik als $(T)$ (tautologie):
\begin{center}
	$\phi$ = $(T) \lor \neg q$ = $T$ \\
\end{center}

Dus $\phi$ is een tautologie omdat Tautologie of ($\lor$) `iets anders' altijd een Tautologie is






\section{Opgave 2}
\textbf{A}

Ik zal eerst de implicaties herschrijven met behulp van de standaardequivalentie $p \rightarrow q$ = $\neg p \lor q$. Dus eerst de linkerkant:
\begin{center}
	$(p \land q) \rightarrow r$ = $\neg(p \land q) \lor r$ = $\neg p \lor \neg p \lor r$ \\
\end{center}

Daardoor word de formule:
\begin{center}
	$\phi$ = $(\neg p \lor \neg q \lor r) \rightarrow (\neg p \land q)$ \\
\end{center}

En dan herschrijf ik beide kanten nogmaals met dezelfde standaardequivalentie. Er komt een $\neg$ (niet) te staan voor het linkerdeel, en de implicatie ($\rightarrow$) wordt een disjunctie ($\lor$):
\begin{center}
	$\phi$ = $\neg(\neg p \lor \neg q \lor r) \lor (\neg p \land q)$ \\
\end{center}

Dan ga ik de negaties binnen de haakjes wegwerken:
\begin{center}
	$\phi$ = $(p \land q \land \neg r) \lor (\neg p \land q)$ \\
\end{center}

En dan heb je de DNF (Disjunctieve Normale Form). Het enige verschil hier is dat ik het nu $\phi_{DNF}$ noem.
\begin{center}
	$\phi_{DNF}$ = $(p \land q \land \neg r) \lor (\neg p \land q)$ \\
\end{center}

\noindent\textbf{B}

\textemdash{}
\vspace{2cm}


\section{Opgave 3}
\textbf{A}

\begin{itemize}
	\item $p$: Je haalt je hand uit de koektrommel.
	\item $q$: Je krijgt een koekje.
\end{itemize}

Formule: $\neg p \rightarrow \neg q$

\noindent\textbf{B}

\begin{itemize}
	\item $p$: Eva geeft een appel aan Adam.
	\item $q$: Adam eet die appel op.
\end{itemize}

Formule: $p \land \neg q$


\end{document}